%%% POST-RECORD v1 BEGIN
% format: hexagon-post-format/1
% record_id: S6_09_PAPERS_I
% version: 1
% date: 2026-11-17
% kind: theorems, computations, corrections, counterexamples, re-derivations
% claim_status: checked results in body; precise open questions in separate final part
% statement_plain: Standard marked family with 0<|w|<216, exact even-shift equivalence, full invariants, and the two analytic boundary obstructions.
% main_statements: fam:src:A.5.1, fam:src:A.5.2, fam:src:C.2.3, fam:src:F.1.1, inv:A3.1, push:src:B.4.33, push:src:D.1.3
% known_before: AC26 construction and cohomology; Engel elliptic-surface description; DGKXv1 standard classification, tangent, Kuranishi, automorphisms and coherent shifts.
% hard_step: actual affine-origin radial defect; exact norm and gradient certificates; complete proof and source comparisons.
% produced_by: GPT-6 Astra (OpenAI, ChatGPT Pro mode, effort Pro), 17 November 2026
% research_runs: GPT-6 Astra (OpenAI, ChatGPT Pro mode), 12-16 November 2026
% working_edition: GPT-5.6 Sol (OpenAI), Codex reasoning effort ultra, 24 August-9 September 2026; a few inputs by GPT-5.6 Sol (OpenAI), ChatGPT Pro mode, 24 August-9 September 2026
% course_continuation: GPT-6 Astra (OpenAI), Codex reasoning effort ultra, October 2026; exact chapter dates retained in the provenance record
% course_prerequisite: GPT-6.1 Sol (OpenAI), reasoning effort ultra, October 2026; full chain-level prerequisite retained with source provenance
% historical_checks: Claude Fable 5.2 (Anthropic), effort high and max, 12-16 November 2026; Claude Opus 5.5 (Anthropic), effort max, 12-16 November 2026
% current_checks: complete proof readings, re-derivations, source-hypothesis checks, exact and interval reruns; exact scopes in concordance
% human_understanding: none claimed; no human has checked the mathematics
% text_origin: AI-produced
% formalization: none
% open_points: literal global coefficient atlas; unconditional normalized norm transcendence
% literature_search: 17 November 2026; original attached literature and exact additional sources; individual locators in inventory
% licence: CC0-1.0
% provenance: companion result inventory, notation dictionary, verification report, scripts, certificates and outputs
%%% POST-RECORD v1 END
% S6_09_PAPERS. Complete proofs and precise open questions are indexed in the inventory.
% AI-produced by GPT-6 Astra (OpenAI, ChatGPT Pro mode, effort Pro), 17 November 2026.
% Text dedicated to the public domain under CC0 1.0.
\documentclass[11pt,a4paper]{article}
\usepackage[margin=25mm,headheight=15pt]{geometry}
\usepackage{fontspec}
\setmainfont{Latin Modern Roman}
\usepackage{amsmath,amssymb,amsthm,amscd,mathtools,mathrsfs}
\usepackage{microtype,booktabs,longtable,array,xcolor,graphicx,enumitem,needspace,fancyhdr}
\usepackage[round,authoryear]{natbib}
\usepackage[unicode,colorlinks=true,linkcolor=blue!40!black,citecolor=blue!40!black,urlcolor=blue!40!black]{hyperref}
\usepackage{xurl}
\usepackage{aliascnt}
\usepackage[nameinlink,noabbrev]{cleveref}
\allowdisplaybreaks[2]
\setcounter{MaxMatrixCols}{32}
\setlength{\emergencystretch}{3em}
\setlength{\parskip}{2pt}
\newtheorem{theorem}{Theorem}[section]
\newaliascnt{proposition}{theorem}
\newtheorem{proposition}[proposition]{Proposition}
\aliascntresetthe{proposition}
\crefname{proposition}{proposition}{propositions}
\newaliascnt{lemma}{theorem}
\newtheorem{lemma}[lemma]{Lemma}
\aliascntresetthe{lemma}
\crefname{lemma}{lemma}{lemmas}
\newaliascnt{corollary}{theorem}
\newtheorem{corollary}[corollary]{Corollary}
\aliascntresetthe{corollary}
\crefname{corollary}{corollary}{corollarys}
\newaliascnt{claim}{theorem}
\newtheorem{claim}[claim]{Claim}
\aliascntresetthe{claim}
\crefname{claim}{claim}{claims}
\newaliascnt{statement}{theorem}
\newtheorem{statement}[statement]{Statement}
\aliascntresetthe{statement}
\crefname{statement}{statement}{statements}
\theoremstyle{definition}
\newaliascnt{definition}{theorem}
\newtheorem{definition}[definition]{Definition}
\aliascntresetthe{definition}
\crefname{definition}{definition}{definitions}
\newaliascnt{remark}{theorem}
\newtheorem{remark}[remark]{Remark}
\aliascntresetthe{remark}
\crefname{remark}{remark}{remarks}
\newaliascnt{example}{theorem}
\newtheorem{example}[example]{Example}
\aliascntresetthe{example}
\crefname{example}{example}{examples}
\newaliascnt{problem}{theorem}
\newtheorem{problem}[problem]{Open problem}
\aliascntresetthe{problem}
\crefname{problem}{open problem}{open problems}
\newaliascnt{question}{theorem}
\newtheorem{question}[question]{Question}
\aliascntresetthe{question}
\crefname{question}{question}{questions}
\newaliascnt{auditrecord}{theorem}
\newtheorem{auditrecord}[auditrecord]{Verification record}
\aliascntresetthe{auditrecord}
\crefname{auditrecord}{verification record}{verification records}
\newaliascnt{test}{theorem}
\newtheorem{test}[test]{Verification}
\aliascntresetthe{test}
\crefname{test}{verification}{verifications}
\newaliascnt{conjecture}{theorem}
\newtheorem{conjecture}[conjecture]{Conjecture}
\aliascntresetthe{conjecture}
\crefname{conjecture}{conjecture}{conjectures}
\newcommand{\C}{\mathbb C}\newcommand{\R}{\mathbb R}\newcommand{\Z}{\mathbb Z}\newcommand{\Q}{\mathbb Q}\newcommand{\N}{\mathbb N}\newcommand{\F}{\mathbb F}
\newcommand{\CC}{\mathbb C}\newcommand{\RR}{\mathbb R}\newcommand{\ZZ}{\mathbb Z}\newcommand{\QQ}{\mathbb Q}\newcommand{\NN}{\mathbb N}
\newcommand{\PP}{\mathbb P}\newcommand{\Pp}{\mathbb P}\newcommand{\Pone}{\mathbb P^1}\newcommand{\HH}{\mathbb H}
\newcommand{\RP}{\mathbb{RP}}\newcommand{\CP}{\mathbb{CP}}
\newcommand{\OO}{\mathcal O}\newcommand{\cO}{\mathcal O}\newcommand{\cM}{\mathcal M}\newcommand{\cF}{\mathcal F}
\newcommand{\h}{\mathfrak h}\newcommand{\hz}{\mathfrak h_z}\newcommand{\Acon}{\mathfrak A}\newcommand{\wt}[1]{\widehat{#1}}
\newcommand{\red}{\mathrm{red}}\newcommand{\tor}{\mathrm{tor}}\newcommand{\dP}{\mathrm{dP}}
\newcommand{\gh}{\widehat\gamma}\newcommand{\uh}{\widehat u}\newcommand{\wh}{\widehat w}\newcommand{\dhat}{\widehat\delta}\renewcommand{\dh}{\widehat\delta}
\newcommand{\dd}{\mathrm d}\newcommand{\ii}{\mathrm i}\newcommand{\eps}{\varepsilon}\newcommand{\partt}{\mathrm{part}}\newcommand{\Bcirc}{B^\circ}\newcommand{\ddbar}{\partial\bar\partial}
\DeclareMathOperator{\Pic}{Pic}\DeclareMathOperator{\Alb}{Alb}\DeclareMathOperator{\Aut}{Aut}\DeclareMathOperator{\NS}{NS}
\DeclareMathOperator{\Hom}{Hom}\DeclareMathOperator{\End}{End}\DeclareMathOperator{\Ext}{Ext}\DeclareMathOperator{\Tor}{Tor}
\DeclareMathOperator{\Spec}{Spec}\DeclareMathOperator{\Proj}{Proj}\DeclareMathOperator{\Sym}{Sym}\DeclareMathOperator{\Diff}{Diff}
\DeclareMathOperator{\GL}{GL}\DeclareMathOperator{\SL}{SL}\DeclareMathOperator{\PSL}{PSL}\DeclareMathOperator{\Sp}{Sp}
\DeclareMathOperator{\im}{im}\DeclareMathOperator{\coker}{coker}\DeclareMathOperator{\rank}{rank}\DeclareMathOperator{\diag}{diag}
\DeclareMathOperator{\ord}{ord}\DeclareMathOperator{\vol}{vol}\DeclareMathOperator{\pr}{pr}\DeclareMathOperator{\id}{id}
\DeclareMathOperator{\divisor}{div}\DeclareMathOperator{\Tr}{Tr}\DeclareMathOperator{\Nm}{Nm}\DeclareMathOperator{\Sha}{Sha}
\DeclareMathOperator{\Imn}{Im}\DeclareMathOperator{\Ren}{Re}\DeclareMathOperator{\trdeg}{tr.deg}
\newcommand{\Uses}[1]{\par\smallskip\noindent\textit{Uses:} #1.\par\smallskip}
\newcommand{\Verified}{}\newcommand{\AuthorSelfCheck}{}\newcommand{\NeedsCheck}{}\newcommand{\Open}{\textbf{Open.}\ }
\pagestyle{fancy}\fancyhf{}\fancyhead[L]{\small The $(3,4,\infty)$ family}\fancyhead[R]{\small Paper I}\fancyfoot[C]{\thepage}
\hypersetup{pdftitle={The (3,4,infinity) family of complex structures on S6: periods, moduli and the ends of the disc},pdfauthor={AI-produced by GPT-6 Astra (OpenAI, ChatGPT Pro mode, effort Pro), 17 November 2026}}

